\subsection{内乘：双代数情形}

设 E 是一个分次双代数（相应地，斜分次双代数）（第 4 号，定义 3）；那么第 6 号和第 7 号的结果可以用来定义右（相应地，左）内乘 \(x \llcorner u \in \mathbf{E}\) 和 \(u \llcorner x \in \mathbf{E}^{*\mathbf{gr}}\) （相应地 \(u \lrcorner x \in \mathbf{E}\) 和 \(x \lrcorner u \in \mathbf{E}^{*\mathbf{gr}}\) )

这对所有 \(x \in \mathbf{E}\) 及所有 \(u \in \mathbf{E}^{*\mathbf{gr}}\) 进行。于是就在 E 上得到一个 (E, E)-双模结构和一个 \((\mathbf{E}^{*\mathbf{gr}}, \mathbf{E}^{*\mathbf{gr}})\)-双模结构。进一步：

命题 9. 设 E 是一个分次双代数（相应地，斜分次双代数）。对 E 中次数为 1 的每个元素 x，被 x 所作的左内乘和右内乘都是代数 E*gr 的导子（相应地，反导子）（§ 10，第 2 号）。

在第 6 号的记号中，对分次双代数（相应地，斜分次双代数）E 中每个次数为 1 的齐次元素 \(x\) ，

\[
c (x) = x \otimes 1 + 1 \otimes x,
\]

这由第 3 号命题 5 及 \(c\) 是次数为 0 的同态这一事实得出。首先假设 \(E\) 是分次双代数。对所有 \(y \in E\) ，根据定义

\[
\langle (u v) \llcorner x, y \rangle = \langle u v, x y \rangle = m ((u \otimes v) (c (x y)))
\]

并且由于 \(c\) 是代数同态， \(c(xy) = c(x)c(y)\) 。设 \(c(y) = \sum_{i} s_i \otimes t_i\) ，其中 \(s_i\) 和 \(t_i\) 属于 \(E\) ；因此

\[
c (x y) = \sum_ {i} (x s _ {i}) \otimes t _ {i} + \sum_ {i} s _ {i} \otimes (x t _ {i}).
\]

。因此 \(\langle (uv) \llcorner x, y \rangle = \sum_{i} u(xs_i)v(t_i) + \sum_{i} u(s_i)v(xt_i)\) 。但我们可以写

\[
\sum_ {i} u (x s _ {i}) v (t _ {i}) = m (((u \llcorner x) \otimes v) (c (y))) = \langle (u \llcorner x) v, y \rangle
\]

并且类似地

\[
\sum_ {i} u (s _ {i}) v (x t _ {i}) = m ((u \otimes (v \llcorner x)) (c (y))) = \langle u (v \llcorner x), y \rangle ,
\]

由此，回到用 \(i(x)\) 记内乘的记号，

\[
(i (x)) (u v) = ((i (x)) (u)) v + u ((i (x)) (v))\tag{64}
\]

，这证明了 \(i(x)\) 是 \(\mathbf{E}^{*}\mathbf{gr}\) 中的导子。

现在假设 \(\mathbf{E}\) 是斜分次双代数，并设 \(u\in \mathbf{E}_p^*\) , \(v\in \mathbf{E}_q^*\) 和 \(y\in \mathbf{E}_r\) ；则我们可以写

\[
c (y) = \sum_ {0 \leqslant j \leqslant r} \left(\sum_ {i} s _ {i j} \otimes t _ {i, r - j}\right)
\]

其中 \(s_{ij}\) 和 \(t_{ij}\) 属于 \(\mathbf{E}_j\) ；则由 \(\mathbf{E}^{\mathfrak{g}} \otimes_{\mathbf{A}} \mathbf{E}\) 中乘积的定义，

\[
c (x y) = c (x) c (y) = \sum_ {0 \leqslant j \leqslant r} \left(\sum_ {i} \left(x s _ {i j}\right) \otimes t _ {i, r - j} + (- 1) ^ {j} \sum_ {i} s _ {i j} \otimes \left(x t _ {i, r - j}\right)\right)
\]

从而这一次

\[
\langle (u v) \llcorner x, y \rangle = \sum_ {0 \leqslant j \leqslant r} \left(\sum_ {i} u \left(x s _ {i j}\right) v \left(t _ {i, r - j}\right) + (- 1) ^ {j} \sum_ {i} u \left(s _ {i j}\right) v \left(x t _ {i, r - j}\right)\right).
\]

同样地， \(\sum_{0\leqslant j\leqslant r}\left(\sum_{i}u(xs_{ij})v(t_{i,r - j})\right) = \langle (u\llcorner x)v,y\rangle\) 。另一方面， \(u(s_{ij}) = 0\) ，除非 \(j = -p\) ，因此我们也可以写成

\[
\sum_ {0 \leqslant j \leqslant r} (- 1) ^ {j} \left(\sum_ {i} u (s _ {i j}) v (x t _ {i, r - j})\right) = (- 1) ^ {p} \langle u (v \llcorner x), y \rangle .
\]

因此我们得出结论：

\[
(i (x)) (u v) = ((i (x)) (u)) v + (- 1) ^ {p} u ((i (x)) (v)),\tag{65}
\]

换言之， \(i(x)\) 是 \(\mathrm{E}^{*\mathrm{gr}}\) 中的反导子。关于 E 中次数为 1 的元素 \(x\) 所作的左内乘的断言，可类似地证明。

注记. (1) 设 E 是 A 上的分次双代数，且 N, N', N'' 是三个分次 A-模。设 m 是从 N × N' 到 N'' 的 A-双线性映射；对于 u \ensuremath{\in} Homgr\_A(E, N) 且 v \ensuremath{\in} Homgr\_A(E, N')，设 u.v 表示 E 到 N'' 的分次同态 m ∘ (u ⊗ v) ∘ c。另一方面，设 i(x) 表示由 x \ensuremath{\in} E 在 A-模 Homgr\_A(E, N), Homgr\_A(E, N') 与 Homgr\_A(E, N'\,') 中所作的（右或左）内乘。那么，若 x 的次数为 1，

\[
(i (x)) (u. v) = ((i (x)) (u)). v + u. ((i (x)) (v))
\]

这对所有 \(u \in \operatorname{Homgr}_{\mathbf{A}}(\mathbf{E}, \mathbf{N})\) 和 \(v \in \operatorname{Homgr}_{\mathbf{A}}(\mathbf{E}, \mathbf{N}')\) 成立。

在相同条件下，若 E 是斜分次双代数且 u 是次数为 p 的齐次元素，则

\[
(i (x)) (u. v) = ((i (x)) (u)). v + (- 1) ^ {p} u. ((i (x)) (v)).
\]

其证明与命题 9 中的相同。

\begin{enumerate}
\def\labelenumi{(\arabic{enumi})}
\setcounter{enumi}{1}
\tightlist
\item
  上述证明中的同一论证更一般地证明了：对所有 \(x \in \mathbf{E}\) ，若 \(c(x) = \sum_{j} x_j' \otimes x_{j'}'\) ，则对所有 \(u, v\) （属于 \(\mathbf{E}^{*\mathrm{gr}}\) ），“莱布尼茨公式”
\end{enumerate}

\[
(i (x)) (u v) = \sum_ {j} (i (x _ {j} ^ {\prime})) (u) \cdot (i (x _ {j} ^ {\prime \prime})) (v)
\]

成立。特别地，对分次双代数 E 的每个本原元素，即满足 \(c(x) = x \otimes 1 + 1 \otimes x\) 的元素， \(i(x)\) 是 \(\mathrm{E}^{*\mathrm{gr}}\) 的导子。

命题 10. 设 E 是一个分次双代数（相应地，斜分次双代数）。对 E*gr 中次数为 1 的任意元素 f，左内乘与右内乘都是代数 E 的导子（相应地，反导子）。

设 \(x \in \mathrm{E}_p, y \in \mathrm{E}_p\) ( \(p \geqslant 1, q \geqslant 1\) ）。根据第 3 号命题 5，我们可以写出

\[
c (x) = x \otimes 1 + \sum_ {1 \leqslant j \leqslant p - 1} \left(\sum_ {i} x _ {i j} ^ {\prime} \otimes x _ {i j, p - j} ^ {\prime \prime}\right) + 1 \otimes x
\]

\[
c (y) = y \otimes 1 + \sum_ {1 \leqslant k \leqslant q - 1} \left(\sum_ {i} y _ {i k} ^ {\prime} \otimes y _ {i, q - k} ^ {\prime \prime}\right) + 1 \otimes y
\]

（其中 \(x_{ij}'\) 和 \(x_{ij}''\) 属于 \(\mathrm{E}_j\) ， \(y_{ik}'\) 和 \(y_{ik}''\) 属于 \(\mathrm{E}_k\) 。若 \(\mathrm{E}\) 是分次双代数，则 \(c(xy) = c(x)c(y)\) 的属于 \(\mathrm{E}_1 \otimes \mathrm{E}\) 的分量等于

\[
\sum_ {i} x _ {i, 1} ^ {\prime} \otimes x _ {i, p - 1} ^ {\prime \prime} y + \sum_ {i} y _ {i, 1} ^ {\prime} \otimes x y _ {i, q - 1} ^ {\prime \prime}
\]

因此根据定义

\[
\begin{array}{r l} (x y) \llcorner f & = \sum_ {i} f (x _ {i, 1} ^ {\prime}) x _ {i, p - 1} ^ {\prime \prime} y + \sum_ {i} f (y _ {i, 1} ^ {\prime}) x y _ {i, q - 1} ^ {\prime \prime} \\ & = (x \llcorner f) y + x (y \llcorner f) \end{array}
\]

而由 \(f\) 所作的右内乘是一个导子。另一方面，若 \(E\) 是斜分次双代数，则 \(c(xy)\) 的属于 \(E_1 \otimes E\) 的分量等于

\[
\sum_ {i} x _ {i, 1} ^ {\prime} \otimes x _ {i, p - 1} ^ {\prime \prime} y + (- 1) ^ {p} \sum_ {i} y _ {i, 1} ^ {\prime} \otimes x y _ {i, q - 1} ^ {\prime \prime}
\]

而这一次我们得到

\[
(x y) \llcorner f = (x \llcorner f) y + (- 1) ^ {p} x (y \llcorner f)
\]

这表明 \(i(f)\) 是一个反导子。对由 \(f\) 所作的左内乘，论证类似。

例子。命题 9 和 10 特别适用于分次双代数 \(\mathsf{S}(\mathbf{M})\) 以及斜分次双代数 \(\bigwedge (\mathbf{M})\) 。由 \(\mathsf{S}(\mathbf{M})\) （相应地 \(\mathsf{S}(\mathbf{M})^{*\mathrm{gr}}\) ）中次数为 1 的元素所作的内乘是彼此可交换的导子，因为 \(\mathsf{S}(\mathbf{M})\) （相应地 \(\mathsf{S}(\mathbf{M})^{*\mathrm{gr}}\) ）是交换的。

类似地，由 \(\bigwedge (\mathbf{M})\) （相应地 \(\bigwedge (\mathbf{M})^{*\mathrm{gr}}\) ）中次数为 1 的元素所作的内乘是平方为零的反导子，因为代数 \(\bigwedge (\mathbf{M})\) （相应地 \(\bigwedge (\mathbf{M})^{*\mathrm{gr}}\) ）中次数为 1 的元素的平方为零。

